\begin{table}[!htbp]
\centering
\begin{threeparttable}
\caption{Distribution of observational school value added}
\label{tab:school-va-distribution}
\small
\setlength{\tabcolsep}{3.5pt}
\begin{tabular}{lcc|cccccc}
\toprule
 & \multicolumn{2}{c|}{Outcome distribution} & \multicolumn{6}{c}{Value-added distribution} \\
\cmidrule(lr){2-3} \cmidrule(lr){4-9}
Outcome & \shortstack{Sample\\Mean} & \shortstack{Sample\\SD} & VA SD & P10 & P25 & P50 & P75 & P90 \\
\midrule
Math & 0.043 & 1.002 & 0.258 & -0.216 & -0.155 & -0.074 & 0.058 & 0.328 \\
Language & 0.049 & 0.989 & 0.188 & -0.202 & -0.135 & -0.030 & 0.095 & 0.238 \\
STEM enrollment & 0.252 & 0.434 & 0.078 & -0.078 & -0.045 & -0.012 & 0.031 & 0.086 \\
Institutional quality & 4.678 & 2.379 & 0.375 & -0.458 & -0.214 & -0.003 & 0.231 & 0.459 \\
\bottomrule
\end{tabular}
\begin{tablenotes}[flushleft]
\footnotesize
\item Notes: The table reports the student-weighted distribution of All-sample adjusted school value added. Schools are weighted by the number of students used to estimate the corresponding value-added measure. Sample Mean and Sample SD refer to the underlying outcome in the outcome-specific value-added regression sample; the remaining columns describe the value-added distribution. Math and language are measured in admission-test standard deviations, STEM enrollment in probability units, and institutional quality in accreditation years.
\end{tablenotes}
\end{threeparttable}
\end{table}
